% GENERATED FILE. Edit canonical lessons, then run npm run generate:books.
% canonical-source-hash: fnv1a64:23615d3da3c23761
% canonical-lessons: RU-C351-poskolku, RU-C351-ved, RU-C351-blagodarya, RU-C351-vsledstvie, RU-C351-sledovatelno

\chapter{Since and thanks to}
\label{ch:ru-since-and-thanks-to}
\hlchaptermodality{351}

\begin{chapteropening}
\textbf{By the end of this chapter:} \emph{I can give a reason with since, after all, thanks to, owing to and therefore, and know what each one takes.}

Complete the last lesson of chapter 351: I can give a reason with since, after all, thanks to, owing to and therefore, and know what each one takes.
\end{chapteropening}

\section[\texorpdfstring{\ru{поскольку} (poskólku)}{поскольку (poskólku)}]{\ru{поскольку} (poskólku) — since, as (giving a reason)}

\label{lesson:RU-C351-poskolku}



\begin{quote}
Before the new one: say the Russian for an effect, then the Russian for as a result.
\end{quote}

\subsection*{You'll want to know: \ru{поскольку}}
\textbf{\ru{поскольку}} (\emph{poskólku}) — \textquotedblleft{}since, as (giving a reason)\textquotedblright{}.

\begin{quote}
\ru{Поскольку} \ru{шёл} \ru{дождь}, \ru{мы} \ru{остались} \ru{дома}.
\end{quote}

\begin{quote}
\emph{Poskólku shyol dozhd, my ostális dóma.}
\end{quote}

\begin{quote}
Since it was raining, we stayed at home.
\end{quote}

It takes an ordinary clause in the indicative, like \emph{\ru{так} \ru{как}}, and sounds a little more formal. It can open the sentence.

\begin{cousinweb}
\textbf{\ru{поскольку}} is built from \emph{\ru{по}} and \emph{\ru{сколько}} \textquotedblleft{}how much\textquotedblright{}, from Proto-Slavic \emph{*koliko} \textquotedblleft{}how much\textquotedblright{} (EtymDB): first \textquotedblleft{}in so far as\textquotedblright{}, then \textquotedblleft{}since\textquotedblright{}.
\end{cousinweb}

\begin{grammarlens}[title={What you've built}]
One more word for telling, explaining and describing at B1.
\end{grammarlens}

\subsection*{Your turn}
\begin{itemize}
\raggedright
  \item \emph{Say it:} \emph{poskólku}
  \item \emph{Say it:} \emph{poskólku}, once more
  \item \emph{Say it:} a sentence of your own with \emph{poskólku}
  \item \emph{Recall:} say the Russian for an effect, then the Russian for as a result, then say \emph{poskólku} again
\end{itemize}

\begin{tcolorbox}[breakable,colback=teal!4,colframe=teal!35!black,title={Before you move on}]
What does \ru{поскольку} mean? (\textquotedblleft{}Since, as (giving a reason)\textquotedblright{}.) Say it once more.
\end{tcolorbox}
\section[\texorpdfstring{\ru{ведь} (ved)}{ведь (ved)}]{\ru{ведь} (ved) — after all, you know}

\label{lesson:RU-C351-ved}



\begin{quote}
Before the new one: say the Russian for as a result, then the Russian for since.
\end{quote}

\subsection*{You'll want to know: \ru{ведь}}
\textbf{\ru{ведь}} (\emph{ved}) — \textquotedblleft{}after all, you know\textquotedblright{}.

\begin{quote}
\ru{Он} \ru{не} \ru{придёт}: \ru{он} \ru{ведь} \ru{болен}.
\end{quote}

\begin{quote}
\emph{On ne pridyót: on ved bólen.}
\end{quote}

\begin{quote}
He will not come: he is ill, after all.
\end{quote}

It gives a reason the listener already knows. It is unstressed and sits near the front of the clause.

\begin{cousinweb}
\textbf{\ru{ведь}} comes from Old East Slavic \emph{vědě} \textquotedblleft{}I know\textquotedblright{} (EtymDB), a form of the old verb \emph{\ru{ведать}} \textquotedblleft{}to know\textquotedblright{}; EtymDB names English \emph{wit} as a cousin. \emph{\ru{Ведь}} still appeals to what both speakers know.
\end{cousinweb}

\begin{grammarlens}[title={What you've built}]
One more word for telling, explaining and describing at B1.
\end{grammarlens}

\subsection*{Your turn}
\begin{itemize}
\raggedright
  \item \emph{Say it:} \emph{ved}
  \item \emph{Say it:} \emph{ved}, once more
  \item \emph{Say it:} a sentence of your own with \emph{ved}
  \item \emph{Recall:} say the Russian for as a result, then the Russian for since, then say \emph{ved} again
\end{itemize}

\begin{tcolorbox}[breakable,colback=teal!4,colframe=teal!35!black,title={Before you move on}]
What does \ru{ведь} mean? (\textquotedblleft{}After all, you know\textquotedblright{}.) Say it once more.
\end{tcolorbox}
\section[\texorpdfstring{\ru{благодаря} (blagodaryá)}{благодаря (blagodaryá)}]{\ru{благодаря} (blagodaryá) — thanks to}

\label{lesson:RU-C351-blagodarya}



\begin{quote}
Before the new one: say the Russian for since, then the Russian for after all.
\end{quote}

\subsection*{You'll want to know: \ru{благодаря}}
\textbf{\ru{благодаря}} (\emph{blagodaryá}) — \textquotedblleft{}thanks to\textquotedblright{}.

\begin{quote}
\ru{Благодаря} \ru{тебе} \ru{я} \ru{сдал} \ru{экзамен}.
\end{quote}

\begin{quote}
\emph{Blagodaryá tebé ya sdal ekzámen.}
\end{quote}

\begin{quote}
Thanks to you I passed the exam.
\end{quote}

It takes the dative and names a good cause: \emph{\ru{благодаря} \ru{помощи}}. For a bad cause Russian says \emph{\ru{из}-\ru{за}} + genitive.

\begin{cousinweb}
\textbf{\ru{благодаря}} is the gerund of \emph{\ru{благодарить}} \textquotedblleft{}to thank\textquotedblright{}, taken from Church Slavonic \emph{blagodariti} (EtymDB), a compound of \emph{\ru{благо}} \textquotedblleft{}good\textquotedblright{} and \emph{\ru{дар}} \textquotedblleft{}gift\textquotedblright{}. The \emph{-\ru{ла}-} is the Church Slavonic shape.
\end{cousinweb}

\begin{grammarlens}[title={What you've built}]
One more word for telling, explaining and describing at B1.
\end{grammarlens}

\subsection*{Your turn}
\begin{itemize}
\raggedright
  \item \emph{Say it:} \emph{blagodaryá}
  \item \emph{Say it:} \emph{blagodaryá}, once more
  \item \emph{Say it:} a sentence of your own with \emph{blagodaryá}
  \item \emph{Recall:} say the Russian for since, then the Russian for after all, then say \emph{blagodaryá} again
\end{itemize}

\begin{tcolorbox}[breakable,colback=teal!4,colframe=teal!35!black,title={Before you move on}]
What does \ru{благодаря} mean? (\textquotedblleft{}Thanks to\textquotedblright{}.) Say it once more.
\end{tcolorbox}
\section[\texorpdfstring{\ru{вследствие} (vslédstviye)}{вследствие (vslédstviye)}]{\ru{вследствие} (vslédstviye) — owing to, in consequence of}

\label{lesson:RU-C351-vsledstvie}



\begin{quote}
Before the new one: say the Russian for after all, then the Russian for thanks to.
\end{quote}

\subsection*{You'll want to know: \ru{вследствие}}
\textbf{\ru{вследствие}} (\emph{vslédstviye}) — \textquotedblleft{}owing to, in consequence of\textquotedblright{}.

\begin{quote}
\ru{Поезд} \ru{опаздывает} \ru{вследствие} \ru{снегопада}.
\end{quote}

\begin{quote}
\emph{Póyezd opázdyvayet vslédstviye snegopáda.}
\end{quote}

\begin{quote}
The train is late owing to the snowfall.
\end{quote}

It takes the genitive. Formal, written style: notices and the news.

\begin{cousinweb}
\textbf{\ru{вследствие}} is fused from \emph{\ru{в}} \textquotedblleft{}in\textquotedblright{} and \emph{\ru{следствие}} \textquotedblleft{}a result\textquotedblright{}, built on \emph{\ru{след}} \textquotedblleft{}a track, a trail\textquotedblright{} (EtymDB): in the wake of something.
\end{cousinweb}

\begin{grammarlens}[title={What you've built}]
One more word for telling, explaining and describing at B1.
\end{grammarlens}

\subsection*{Your turn}
\begin{itemize}
\raggedright
  \item \emph{Say it:} \emph{vslédstviye}
  \item \emph{Say it:} \emph{vslédstviye}, once more
  \item \emph{Say it:} a sentence of your own with \emph{vslédstviye}
  \item \emph{Recall:} say the Russian for after all, then the Russian for thanks to, then say \emph{vslédstviye} again
\end{itemize}

\begin{tcolorbox}[breakable,colback=teal!4,colframe=teal!35!black,title={Before you move on}]
What does \ru{вследствие} mean? (\textquotedblleft{}Owing to, in consequence of\textquotedblright{}.) Say it once more.
\end{tcolorbox}
\section[\texorpdfstring{\ru{следовательно} (slédovatelno)}{следовательно (slédovatelno)}]{\ru{следовательно} (slédovatelno) — therefore, consequently}

\label{lesson:RU-C351-sledovatelno}



\begin{quote}
Before the new one: say the Russian for thanks to, then the Russian for owing to.
\end{quote}

\subsection*{You'll want to know: \ru{следовательно}}
\textbf{\ru{следовательно}} (\emph{slédovatelno}) — \textquotedblleft{}therefore, consequently\textquotedblright{}.

\begin{quote}
\ru{Он} \ru{не} \ru{звонил}, \ru{следовательно}, \ru{всё} \ru{в} \ru{порядке}.
\end{quote}

\begin{quote}
\emph{On ne zvoníl, slédovatelno, vsyo v poryádke.}
\end{quote}

\begin{quote}
He has not called, so everything is all right.
\end{quote}

A formal link between a premise and a conclusion, set off by commas. In conversation people say \emph{\ru{значит}}.

\begin{cousinweb}
\textbf{\ru{следовательно}} is built on \emph{\ru{следовать}} \textquotedblleft{}to follow\textquotedblright{}, from \emph{\ru{след}} \textquotedblleft{}a track, a trail\textquotedblright{} (EtymDB): what follows from what came before.
\end{cousinweb}

\begin{grammarlens}[title={What you've built}]
One more word for telling, explaining and describing at B1.
\end{grammarlens}

\subsection*{Your turn}
\begin{itemize}
\raggedright
  \item \emph{Say it:} \emph{slédovatelno}
  \item \emph{Say it:} \emph{slédovatelno}, once more
  \item \emph{Say it:} a sentence of your own with \emph{slédovatelno}
  \item \emph{Recall:} say the Russian for thanks to, then the Russian for owing to, then say \emph{slédovatelno} again
\end{itemize}

\begin{tcolorbox}[breakable,colback=teal!4,colframe=teal!35!black,title={Before you move on}]
What does \ru{следовательно} mean? (\textquotedblleft{}Therefore, consequently\textquotedblright{}.) Say it once more.
\end{tcolorbox}
